\clearpage\thispagestyle{empty}
\pdfbookmark[0]{Quellen zu Vesta (24. Dezember 1812)}{vesta1812dec24-sources}
\begin{center}\textbf{Quellen zu Vesta (24. Dezember 1812)}\end{center}
\noindent Diese KI-Übertragung ins heutige Deutsch folgt dem Bericht vom 24. Dezember 1812 in den \href{https://gdz.sub.uni-goettingen.de/id/PPN236007467}{\textit{Werken}, Band~VI} von Gauß (Göttingen, 1874), S.~361--364. Der Erstdruck steht unter der Überschrift Göttingen in der Zeitschrift \textit{Göttingische gelehrte Anzeigen}, Stück~205, S.~\href{https://api.digitale-sammlungen.de/iiif/image/v2/bsb10538682_01007/full/1600,/0/default.jpg}{2041}, \href{https://api.digitale-sammlungen.de/iiif/image/v2/bsb10538682_01008/full/1600,/0/default.jpg}{2042}, \href{https://api.digitale-sammlungen.de/iiif/image/v2/bsb10538682_01009/full/1600,/0/default.jpg}{2043}, \href{https://api.digitale-sammlungen.de/iiif/image/v2/bsb10538682_01010/full/1600,/0/default.jpg}{2044}, \href{https://api.digitale-sammlungen.de/iiif/image/v2/bsb10538682_01011/full/1600,/0/default.jpg}{2045}. Die Links führen zum digitalisierten Exemplar der Bayerischen Staatsbibliothek. Zum Bericht gehört der abschließende Abschnitt über Nicolais Arbeit zum zweiten Kometen von 1811.
\par\smallskip\noindent Beide Ausgaben lassen den Tag in der Zeile zur Oktoberopposition offen. Bei Enkes Elementen auf S.~2044 endet die Länge des Perihels im Erstdruck mit $31{,}1$. Im Sinusausdruck zur Exzentrizität stehen dort Grad-, Minuten- und Sekundenzeichen. Die Werkausgabe gibt die Länge als $249^\circ\;38^{\prime}\;31.4$ und das Sinusargument ohne diese Zeichen an. KI-Übertragung und KI-Satz.
\clearpage
